\scanpage{93}
\mspage{\handwritten{1}}
% Only the bottoms of strokes survive at the top right of this copy.
\margin{\illegible{3}}

\noindent Dear Daniel,\\
\margin{\handwritten{1.}}Last year Ronnie Brown from Bangor sent me a heap of reprints and preprints by him and a group of friends\handwritten{,} on various foundational matters of homotopical algebra. I did not really dig through any of this, as I kind of lost contact with the technicalities of this kind (I was never too \unsure{familiar} with the homotopy techniques \inserted{\handwritten{anyhow,}} I confess) - but this reminded me of a few letters I had exchanged with Larry Breen in\unsure{.}1975, where I had developped an outline of a \inserted{program for a} kind of ``topological algebra'', viewed as a synthesis of \deleted{\unsure{ho}}\unsure{homotopical} and homological algebra, with special emphasis on topoi - most of the basic intuitions in this program arising from various \unsure{backgrounds|background} in algebraic geometry. \deleted{(}Some of those intuitions we discussed\handwritten{,} I believe\handwritten{,} at IHES eight or nine years before, at a time when you had just written up your nice ideas on axiomatic homotopical algebra, published since in Spinger's Lec\unsure{t}ure Notes. I write you under the assumption that you have not entirely lost interest for those foundational questions you were looking at more than fifteen years ago. One thing which strikes me, is that (as far as I know) there has not been any substantial progress since - it l\deleted{\unsure{a}}ooks to me that \deleted{\unsure{adequate}} understanding of the basic structures underlying homotopy theory, or even homological algebra only, is still \unsure{lacking} - \unsure{probably} because the few people who have a wide enough background and perspective enabling them to feel the main questions, are devoting their energies to things which seem more directly rewarding. Maybe \deleted{there} even a wind of disrepute for any foundational matters whatever is blowing nowadays ! In this respect, what seems to me even more striking than the lack of proper foundations \handwritten{f}or homological and homotopical algebra, is the absence I dare\unsure{-}say\unsure{,} of proper foundations for topology itself ! I am thinking here \unsure{mainly} of the development of a context of ``tame'' topology, which (I am \unsure{convinced)} would have on the everyday technique of geometric topology (\deleted{in} I use this expression in contrast to the topology of use for analysts) a comparable impact or even a greater one, than the introduction of the point of view of schemes had on algebraic geometry. The \inserted{psychological} drawback here I believe is not anything like messyness, as for homological and homotopical algebra\inserted{\handwritten{\unsure{(or for schemes)}}}, but \unsure{merely} the inrooted intertia which \deleted{\unsure{goes against}} prevents \deleted{\unsure{everyone}} us so stubbornly from looking innocently, with fresh eyes, upon things, without being dulled and emprisoned by standing habits of thought, going with a familiar context - \uline{too} familiar a context ! The task of working out the foundations of tame topology\handwritten{,} and a \deleted{theory} corresponding structure theory for ``stratified (tame) spaces\handwritten{''}, seems to me a lot more urgent and exciting still \unsure{than}
% archive sheet number: 92
